\chapter{Testing}
\label{chap:testing}

The platform's behaviour splits into two classes that demand different verification strategies. \emph{Architectural correctness} --- the contract between client and server, the integrity of the credit ledger, the correctness of authentication, and the shape of the data crossing the multimodal abstraction --- is amenable to automated testing and is enforced on every pull request. The \emph{qualitative output} of the AI providers is stochastic and is evaluated offline against fixed reference sets and rubrics. The two suites taken together are the platform's quality bar.

\section{Backend integration tests}
\label{sec:backend-tests}

The backend test project mirrors the production code's feature folders, with one xUnit test class per endpoint area. The harness brings up a fresh PostgreSQL instance with the pgvector extension enabled, applies the schema migrations, starts the ASP.NET Core application in-process, and drives it through HTTP. The database is deliberately not mocked: pgvector behaviour and EF Core migration semantics are part of what the suite is meant to exercise.

External AI providers are stubbed at the seam. The \texttt{IImageProvider} interface is replaced by a \texttt{MockImageProvider} that returns a deterministic placeholder image, and the recommendation engine's language-model calls are routed through a stub kernel that returns canned outputs for canned inputs. The substitution makes the suite reproducible and free to run, while leaving prompt construction, embedding storage, ledger debits and job scheduling fully exercised.

\section{End-to-end tests with Playwright}
\label{sec:client-tests}

The web build of the React Native client is exercised by Playwright against the development environment, which uses a separate database, a separate blob account and the mock image provider so that no real generation cost is incurred during a test run. The suite walks the user-visible journeys end to end: account creation through the three-stage onboarding flow, product creation with background removal and search-by-photo retrieval, generation of one asset of each family (announcement, catalogue, wallpaper) under both fully generative and hybrid modes, A4 print export, and the daily-recommendation handoff into the asset-generation pipeline. The credit ledger is observed throughout, so that the wallet balance after any sequence of operations always matches the sum of the corresponding ledger entries.

Because the React Native web export is a real DOM tree and the iOS/Android builds reuse the same components, a passing Playwright run is strong evidence that the route tree, the state-context layer and the API client are correct on all three surfaces. Native-only behaviour (haptics, OS deep links) is verified manually before each release.

\section{LLM output evaluation}
\label{sec:llm-eval}

The recommendation engine's writer stage produces stochastic text and cannot be verified by equality assertions. It is instead evaluated offline against a frozen reference set of $120$ $(\text{shop}, \text{day}, \text{signal})$ tuples that span urban and rural shop profiles, weak and strong context signals, and both supported languages. Three complementary checks are run on every model or prompt change:

\begin{itemize}
    \item \emph{Schema and format gate.} Every output must parse as the declared JSON schema, populate every required field and respect the fixed enumerations for \texttt{type} and \texttt{tone}. Schema failures fail the suite.
    \item \emph{Forbidden-phrase regression.} A curated denylist of the boilerplate marketing tropes the writer prompt is built to discourage (``discover our premium selection'', ``don't miss out'', and their Romanian equivalents) is matched against the \texttt{title}, \texttt{body} and \texttt{suggestedPost} fields. The suite tracks the per-100-ideas rate of denylist hits and gates merges on it not regressing.
    \item \emph{LLM-as-judge rubric scoring.} A stronger reference model (\emph{Claude} or \emph{Gemini}, alternated to control for judge bias) scores each output on a $1$--$5$ rubric covering tonal fit, specificity to the trigger signal, and call-to-action concreteness, following the LLM-as-judge protocol of Zheng et al.~\cite{zheng2023llmjudge}. The suite reports mean scores per axis with bootstrap confidence intervals and flags regressions of more than $0.2$ rubric points.
\end{itemize}

\noindent A complementary human-rated pass is run before each release. Two independent reviewers perform pairwise comparisons of the platform's output against a baseline (the un-adapted seven-billion-parameter writer) on a $30$-tuple stratified sub-sample, and the win-rate is recorded. Image generation is evaluated only manually: a small canonical prompt set is run against each live provider and the resulting assets are inspected by a human operator for product fidelity (in hybrid mode) and tonal fit.

\section{Continuous integration and the merge gate}
\label{sec:ci-merge-gate}

The build pipeline is hosted on GitHub Actions. On every pull request it runs the backend integration suite against an ephemeral PostgreSQL service container with pgvector enabled, runs the Playwright suite against the development environment, and re-runs the LLM-output schema gate and forbidden-phrase regression on the frozen reference set. All three must pass before merge into \texttt{master}; the merge gate is the platform's single quality bar with no manual override. Production deployment proceeds from the merged commit through the same pipeline and runs no further test gating. Table~\ref{tab:capability-test-layer} maps each functional capability area to the test layer that protects it.

\begin{table}[!ht]
\centering
\small
\renewcommand{\arraystretch}{1.3}
\begin{tabularx}{\textwidth}{|>{\raggedright\arraybackslash}p{4.0cm}|Y|Y|Y|}
\hline
\textbf{Capability area} & \textbf{Backend integration} & \textbf{Playwright E2E} & \textbf{LLM evaluation} \\
\hline
Identity and onboarding   & yes & yes & --- \\
\hline
Product library           & yes & yes & --- \\
\hline
Reference search          & yes & yes & --- \\
\hline
Asset generation          & yes (mock provider) & yes (mock provider) & manual on live providers \\
\hline
Daily recommendations     & yes (stub kernel) & yes (read path only) & schema, denylist, LLM-as-judge \\
\hline
Library and print         & yes & yes & --- \\
\hline
Billing                   & yes & yes & --- \\
\hline
Administration            & yes & yes & --- \\
\hline
\end{tabularx}
\caption{Mapping from functional capability area to the test layer that protects it.}
\label{tab:capability-test-layer}
\end{table}
